%======================================================================
% CardamomOT model learning: project note
%======================================================================
\documentclass[a4paper, 11pt]{article}
\usepackage[utf8]{inputenc}
\usepackage[T1]{fontenc}
\usepackage{amsmath, amssymb, amsthm}
\usepackage[margin=1in]{geometry}
\usepackage{enumitem}
\usepackage{booktabs}
\usepackage{hyperref}
\hypersetup{colorlinks=true, linkcolor=blue, citecolor=blue, urlcolor=blue}

\DeclareMathOperator*{\argmin}{arg\,min}
\DeclareMathOperator*{\argmax}{arg\,max}
\newcommand{\E}{\mathbb{E}}
\newcommand{\kon}{k_{\mathrm{on}}^{\theta}}
\newcommand{\koni}{k_{\mathrm{on},i}^{\theta}}
\newcommand{\CE}{\mathrm{CE}}
\newcommand{\F}{\mathcal{F}}
\newcommand{\tofill}[1]{\par\smallskip\noindent\emph{[To be completed: #1]}\par\smallskip}

\theoremstyle{plain}
\newtheorem{prop}{Proposition}
\theoremstyle{remark}
\newtheorem{rem}{Remark}

\title{CardamomOT-ML\\[0.3em] \large Learning the regulation function of a mechanistic gene regulatory network\\ from single-cell RNA time series}
\author{Project note}
\date{\today}

\begin{document}
\maketitle

%======================================================================
\section*{Project note}
%======================================================================

\paragraph{Objective.} CardamomOT~\cite{mauge2026cardamomot} infers a gene regulatory network (GRN), hidden protein trajectories and kinetic rates from scRNA-seq time series, with a burst frequency that is a sigmoid of a linear combination of the proteins. This shape is an assumption, inherited from Harissa~\cite{herbachinferring2017}, not a result. The article will make the regulation function a quantity learned \emph{gene by gene} during the inference: each target gene chooses, among a family of candidate functions, the one that best explains its basins along the reconstructed trajectories, while the network stays a single matrix $\theta$ whose support and signs keep their meaning.

\paragraph{The thread.} With bimodal burst frequencies, a cell only reveals whether a gene is in its low or high mode; the shape of the regulation function between the two modes is invisible at a steady state. It becomes visible along \emph{transients}: proteins relax slowly ($1/d_1$), so the reconstructed trajectories cross intermediate protein levels, and the time at which a target switches, relative to the protein level of its regulators, depends on the shape of the function (Section~\ref{sec:identifiability}). The mechanistic trajectories of CardamomOT, which time the switches of every gene, are therefore exactly the object on which the regulation function can be learned. The family is selected as a discrete parameter of the EM, by the held-out likelihood of the counts (Section~\ref{sec:selection}).

\paragraph{What already exists in the code.} The inference loop is an EM algorithm (Section~\ref{sec:em}). The counts of each gene are read by a continuous model of its burst frequency, so that a cell caught during a switch is read at an intermediate level; the basin update and the network fit maximise the same likelihood; the network is fitted on the expected levels; the trajectories are forced by a mass-constrained assignment of the same costs; and the regulation function is compared with the targets along the memory of the mRNA (Section~\ref{sec:memory}). The regulation function is computed by a single module, \texttt{CardamomOT/inference/sigma.py}, for the whole inference and every simulator (numba kernels, and a torch mirror for the free-running fit of the kinetic rates). Its families are given by a specification per gene (transformation $\varphi$ of the regulators and link). The network fit supports any $\varphi$ with the logistic link. The selection of Section~\ref{sec:selection} is implemented in \texttt{fit\_network} (\texttt{model\_learning\_grn}), with the candidate families \texttt{model\_learning\_phi = ['identity', 'herbach']} (approximate Herbach form, $m \in \{1, 2, 4\}$; the Hill type on request) and \texttt{model\_learning\_link = ['logistic']}; the proteins of the cells can be given as observed (\texttt{Data/data\_prot.h5ad}, Section~\ref{sec:benchmark}); the chromatin opening function of the multiome extension has its own, independent families (\texttt{model\_learning\_chrom}). The benchmark of Section~\ref{sec:benchmark} (\texttt{projet\_modellearning/}) simulates the networks of the CardamomOT article with sigmoid and Herbach targets, with the model of the inference.

\paragraph{Where we stand.} Earlier benchmarks combined every difficulty at once (explicit mRNA with memory, counts, reconstructed proteins): the families were recovered at $0.55$--$0.68$ only, also on the true proteins, while noise-free targets recovered them all. The benchmark is now built in levels (Section~\ref{sec:benchmark}): data simulated with the model of the inference (proteins only, counts drawn from the negative binomial of the burst frequency), the proteins first observed, then hidden.

\paragraph{Next steps.}
\begin{enumerate}[nosep]
\item Levels 0, 1 and 2 of the benchmark on the four networks (Section~\ref{sec:results}).
\item Then the difficulties one at a time: mRNA with memory (Section~\ref{sec:memory}), partially observed proteins.
\item Candidate families whose differences are not confined to the tails of $\sigma$ (Richards link, Section~\ref{sec:richards}), co-regulation terms (Section~\ref{sec:coreg}).
\item Validation of the selected families in the free-running fit of the kinetic rates (generative criterion), then real data (CITE-seq for observed proteins).
\end{enumerate}

%======================================================================
\section{Introduction}\label{sec:intro}
%======================================================================

Mechanistic GRN models fix the form of the regulation function before the inference: Hill functions in ODE models, sums of sigmoids in neural-network-inspired models, logistic links in the bursty models of single-cell data~\cite{herbachinferring2017, mauge2026cardamomot}. The form matters: it decides how a target responds to an intermediate level of its regulators, hence the timing of switches, the existence of bistability, and the interpretation of an interaction strength. The thermodynamic models of transcription~\cite{shea1985lambda, bintu2005transcriptional} show that the input function of a promoter depends on how transcription factors bind and interact, and differs from gene to gene.

Learning the form of a dynamical model from data is the object of model selection and of sparse identification of nonlinear dynamics~\cite{brunton2016sindy, mangan2016inferring}, in which the vector field is expanded on a library of candidate terms. These approaches require observed trajectories. Single-cell RNA time series give snapshots of independent cells, with unobserved proteins, and the regulation function acts on bursty transcription. CardamomOT reconstructs the hidden protein trajectories by mechanistic optimal transport; the regulation function can therefore be learned on them, at the level of each target gene, in the same alternation as the network.

%======================================================================
\section{Background: CardamomOT}\label{sec:background}
%======================================================================

The promoter of each gene $i$ produces mRNA in bursts at frequency $\koni(P)$, and the proteins are translated from the mRNA~\cite{herbachinferring2017}:
\begin{equation}\label{eq:pdmp}
M_i \xrightarrow{\ \koni(P)\ } M_i + \mathcal{E}\Bigl(\frac{k_{\mathrm{off},i}}{s_{0,i}}\Bigr), \qquad
\frac{dM_i}{dt} = -d_{0,i}M_i, \qquad \frac{dP_i}{dt} = s_{1,i}M_i - d_{1,i}P_i,
\end{equation}
with
\begin{equation}\label{eq:kon}
\koni(P) = k_{0,i} + (k_{1,i} - k_{0,i})\,\sigma_i(P), \qquad \sigma_i(P) = \sigma\Bigl(\beta_i + \sum_j \theta_{ij}P_j\Bigr),
\end{equation}
$\sigma$ the logistic function and $s_1$ chosen so that $P_i \in [0,1]$ (high mode at 1). The negative binomial mixture of each gene gives the burst parameters and the posterior probabilities of its basins in each cell. The loop then alternates between (1) the coupling of consecutive timepoints by entropic OT with a mechanistic bridge cost, which reconstructs the protein trajectories and the switching time of each gene, (2) the network fit, a cross-entropy between $\sigma_i$ and the position $\tau_i \in [0,1]$ of the basin of the target between $k_0$ and $k_1$, with a sparsity penalty, and (3) the update of the basins, which balances the mixture likelihood and the consistency with the network under constraints on the proportions of the modes. In the published version, steps (2) and (3) use hard basins, and the network term of step (3) is a heuristic proximity between the predicted and the modal burst frequencies. The next section turns the loop into an EM algorithm, which is the ground on which the regulation function can be learned.

%======================================================================
\section{CardamomOT as an EM algorithm}\label{sec:em}
%======================================================================

\paragraph{Observation model.} The counts of gene $i$ follow a negative binomial mixture whose components are burst frequencies of the gene. With a few modes $\mathbf{k}_{0,i} < \dots < \mathbf{k}_{M_i-1,i}$, the mixture is fitted by a classification EM under mode masses per timepoint, then refined by a soft EM with the same masses (\texttt{max\_iter\_softEM}), which removes the truncation bias of the classification EM. The mixture is then the continuous model of the counts:
\begin{equation}\label{eq:grid}
x_{ni} \sim \sum_{l=1}^{L} w_{l,i}(t_n)\,\mathrm{NB}\bigl(x;\ k_{0,i} + (k_{1,i} - k_{0,i})\,\sigma_l,\ c_i\bigr), \qquad \sigma_l = \frac{l-1}{L-1},
\end{equation}
at the depth of the cell, with $L$ levels chosen per gene between $2$ and \texttt{max\_grid\_levels} by a BIC or ICL on the soft likelihood of the refined grids (\texttt{mixture\_selection\_criterion}), weights per timepoint under a Dirichlet prior that favours the extreme levels, and $(k_{0,i}, k_{1,i}, c_i)$ and the weights fitted by EM from the extremes of the mixture with modes. Its levels are the modes of the whole inference, $\mathbf{k}_{l,i} = k_{0,i} + (k_{1,i} - k_{0,i})\sigma_l$: a cell caught during a switch is read at an intermediate level rather than at an extreme mode, and the expected level of a cell under the mixture alone, $\E[\sigma \mid x_{ni}]$, is its position between $k_{0,i}$ and $k_{1,i}$ given its counts.

\paragraph{Latent variables and likelihood.} The basin $z_{ni}$ (mode or level) of gene $i$ in cell $n$ is latent; given the proteins of the trajectory reaching $n$, the complete likelihood of the mRNA $x_{ni}$ and of the basin is
\begin{equation}\label{eq:complete}
p(x_{ni}, z_{ni} = m \mid P_n) \;\propto\; \mathrm{NB}\bigl(x_{ni};\, \mathbf{k}_{m,i}\bigr)\, \exp\bigl(-\kappa_i\, \mathrm{KL}(\tau_{m,i} \,\|\, \bar\sigma_i(n))\bigr),
\end{equation}
with $\tau_{m,i} = (\mathbf{k}_{m,i} - k_{0,i})/(k_{1,i} - k_{0,i})$ the position of mode $m$ between the extreme modes, $\mathrm{KL}(\tau\,\|\,\sigma) = \tau\log(\tau/\sigma) + (1-\tau)\log((1-\tau)/(1-\sigma))$, and $\bar\sigma_i(n)$ the regulation function along the memory of the mRNA of the cell (Section~\ref{sec:memory}; $\bar\sigma_i(n) = \sigma_i(P_n)$ without it). With two modes, $\kappa_i = 1$ and the second factor is $\bar\sigma_i(n)$ or $1 - \bar\sigma_i(n)$: the network gives the probability of the high mode, as in the regulated two-state promoter. With the levels of~\eqref{eq:grid}, $\exp(-\mathrm{KL}(\sigma_l \,\|\, \bar\sigma))$ alone is nearly flat over the levels, of width $\sim\bar\sigma(1 - \bar\sigma)$; the precision $\kappa_i$ sets how sharply the network places a cell at the level $\bar\sigma_i(n)$. It is set per gene at each update by moment matching on the cells reached by the trajectories: the width $\bar\sigma(1-\bar\sigma)/\kappa_i$ of the prior matches the mean second moment $\E_q[(\tau_{z,i} - \bar\sigma_i)^2]$ of the levels under the posteriors, clipped to $[1, \kappa_{\max}]$ (\texttt{grid\_network\_precision\_max}).

\paragraph{E-step: posteriors and basins.} At each timepoint, the cells (uniform mass) are coupled to the modes (masses given by the mixture), with $\pi_{ni}$ the posterior of the mixture alone. (i) The posteriors are the rows of the entropic OT of regularisation $1$ with the cost $-\log\pi_{ni}(m) + \kappa_i\,\mathrm{KL}(\tau_{m,i}\,\|\,\bar\sigma_i(n))$: $q_{ni}(m) \propto \pi_{ni}(m)\, e^{-\kappa_i \mathrm{KL}(\tau_{m,i}\|\bar\sigma_i(n))}\, a_{m,i}$, with one factor $a_{m,i}$ per mode so that the masses match, i.e.\ the posteriors of~\eqref{eq:complete} under the mass constraint. (ii) The basins are the exact OT (EMD) with the cost
\begin{equation}\label{eq:costs}
C_{nm} = -\log\pi_{ni}(m) + (1 - \omega_k)\,\kappa_i\, w_i\, \mathrm{KL}(\tau_{m,i} \,\|\, \bar\sigma_i(n)),
\end{equation}
with $\omega_k = \max(0.96^{k-1}, 0.1)$ at iteration $k$ and $w_i = 1 - k_{0,i}/k_{1,i}$: one basin per cell, with exactly the prescribed masses. The network enters the hard assignment slowly, and weakly for genes whose modes are weakly separated, because a label set by the network is then explained perfectly by it (the classification EM locks in); the mass constraint protects the proportions of the modes against an early network, which can only reassign basins between cells. The posteriors need no such protection: a cell whose counts do not discriminate the modes gets the network prediction as target ($q$ close to the prior), and its cross-entropy gradient vanishes at the current network, so that it slows the updates without orienting them; a bad network is corrected by the informative cells, and its precision $\kappa_i$ falls with their residuals. A weaker network term in (i) would instead give such cells the mean of their timepoint as target, which pulls every network towards a constant ($\theta$ attenuated by the reliability of the counts). The posteriors (i) give the targets of the network fit, the basins (ii) force the trajectories. The network term is the log-likelihood~\eqref{eq:complete}, instead of the proximity of burst frequencies of the published version.

\paragraph{M-step: network fit on the expected levels.} The network is fitted on the expected positions $\bar\tau_{ni} = \E_{q}[\tau_{z_{ni},i}]$ (\texttt{one\_hot\_probas = False}; with \texttt{True}, on the basins (ii), a classification EM). The EM is started from the classification EM: the targets are $\lambda_k\tau_{z^{\mathrm{EMD}}_{ni},i} + (1 - \lambda_k)\bar\tau_{ni}$, with $\lambda_k$ decreasing linearly from $1$ at the first iteration to $0$ at \texttt{min\_n\_loops}. At the first iteration there is no network yet, and the expected positions of the cells whose counts do not discriminate the modes would be the mean of their timepoint: a flat network, from which the EM leaves slowly.

\begin{prop}\label{prop:mstep}
$\sum_n \CE\bigl(\bar\sigma_i(n), \bar\tau_{ni}\bigr) = \sum_n \E_{q}\bigl[\mathrm{KL}(\tau_{z_{ni},i} \,\|\, \bar\sigma_i(n))\bigr] + C$, with $C$ independent of $(\beta_i, \theta_{i\cdot})$. The network fit on the expected levels minimises the expected complete negative log-likelihood~\eqref{eq:complete}: it is the M-step of the EM.
\end{prop}

\begin{proof}
The cross-entropy relative to its target, $\CE(\sigma, \tau) = -\tau\log\sigma - (1-\tau)\log(1-\sigma) - H(\tau)$, is affine in $\tau$ up to the entropy $H(\tau)$, which does not depend on the network; the expectation over $q$ commutes with the affine part.
\end{proof}

The E-step and the M-step thus maximise the same likelihood, and the expected levels give the M-step of the full EM, with the variance of a draw removed (Rao--Blackwellisation). The factors $\kappa_i$ and $(1 - \omega_k)w_i$ weigh the network against the mixture in the E-step; the network fit, gene by gene, minimises the cross-entropy of each target without them.

\paragraph{Trajectories: forcing by the mass-constrained basins.} The trajectories need, for each cell, the burst frequency that forces the proteins along the bridges and decides the switching times. They are forced by the basins (ii), recomputed at every iteration. With ordered levels, the EMD of the costs~\eqref{eq:costs} is a rank-preserving assignment: at each timepoint, the forcing frequencies have exactly the histogram of the mixture, the cells keep the order given by their posteriors and the network, and the assignment is deterministic. These are the triple-goal estimates of hierarchical models~\cite{shen1998triple}, whose histogram and ranks are right, unlike those of the expectations (contracted histogram) or of draws (noisy ranks). The expected levels are the right targets of the network fit (Proposition~\ref{prop:mstep}), but forcing the trajectories with them would contract the distribution of the burst frequencies ($\mathrm{Var}\,\E[\sigma\mid x] < \mathrm{Var}\,\sigma$): a cell with an uncertain basin would be forced at an intermediate frequency for a lack of information rather than because the network says so, which flattens the trajectories and the next network. A stochastic approximation~\cite{delyon1999saem}, which averages draws of the basins along the iterations, converges to the same contracted expectation, and before converging its average carries the posteriors of past iterations, a lag that mixes the current state of the loop with its history. The basins (ii) have neither the contraction nor the lag; with the levels of~\eqref{eq:grid}, a cell in a transient is forced at an intermediate frequency where the counts and the network put it. The initial proteins of the trajectories are the protein levels of the forced modes, per sample.

\paragraph{Free-running fit of the kinetic rates.} Its loss compares distributions (Sinkhorn divergences between simulated and observed burst frequencies), so that its targets must have the right distribution, not the right mean: they are levels drawn from the posteriors $q$ at each minibatch. For genes with two modes, the simulated burst frequency is itself a Bernoulli draw of $\sigma$ between $k_0$ and $k_1$ (straight-through gradient), so that $\sigma$ keeps its meaning of a probability of the high mode.

\paragraph{Stopping criterion.} The error of the loop is the network term of the OT between trajectories and cells, the cross-entropy between $\bar\sigma$ and the targets weighted by $w_i$ (and the flow-matching term when it is used): the transport, the basin update and the stopping criterion measure the same quantity.

\paragraph{Test set.} The basins of the held-out cells come from the same two couplings, the network term being replaced by a classifier calibrated on how the training settled its cells: the EMD basins force their trajectories, as in the training, and the entropic coupling gives their posteriors and expected levels.

\paragraph{Why it matters for model learning.} (i) A single likelihood~\eqref{eq:complete} is maximised by both steps, for any regulation function: comparing two families is comparing two likelihoods, and a new family only has to provide its own $\sigma_i$. (ii) Soft posteriors keep the uncertainty of the basins: with hard labels and a network term, the labels of undecided cells are set by the current regulation function, which then explains them perfectly; the expected levels keep it, and a target changes family when the informative cells favour another one (Section~\ref{sec:selection}). (iii) The continuous model and the forcing by the basins let the trajectories carry intermediate burst frequencies exactly where the counts and the network put them, which is where the shape of the regulation function between the two modes is expressed.

%======================================================================
\section{Targets along the memory of the mRNA}\label{sec:memory}
%======================================================================

The basin of a cell is read on its mRNA, which integrates the burst frequency over its lifetime: by~\eqref{eq:pdmp}, the mean mRNA at time $t$ is proportional to $\int_0^\infty \koni(P(t-u))\, d_{0,i}e^{-d_{0,i}u}\,du$. The target of a cell is therefore compared with the regulation function averaged along this memory,
\begin{equation}\label{eq:memory}
\bar\sigma_i(n) = \sum_{q=0}^{Q} w_{q,i}\, \sigma_i\bigl(P_n(t_n - h_q)\bigr), \qquad h_q = \frac{q}{Q}H, \quad H = \min\Bigl(\Delta t, \frac{5}{\min_j d_{0,j}}\Bigr),
\end{equation}
with $P_n(s)$ the proteins of the bridge reaching cell $n$, $\Delta t$ the interval between timepoints, and $w_{q,i}$ the trapezoidal weights of $d_{0,i}e^{-d_{0,i}h}$, the mass before the window being put on the earliest node. With a single node ($Q = 1$, \texttt{late\_mrna\_nodes = 1}), $\bar\sigma_i(n) = \sigma_i(P_n(t_n - 1/d_0))$, the mean lifetime of the mRNA (median over the genes). The regulation function itself stays instantaneous: the proteins drive transcription, and only its read-out is delayed. There is no circularity, since the proteins are built from the modes read on the mRNA, which is the input of translation.

The average~\eqref{eq:memory} (\texttt{targets\_accounting\_late\_mrna}, $Q$ = \texttt{late\_mrna\_nodes}) is used in the network fit, the basin update (with $\bar\sigma$ averaged over the states that reach each cell), the stopping criterion and the score of the selection. The OT between trajectories and cells and the switching times of the forcing, which evaluate $\sigma$ for every pair of states, use the single node. The flow-matching nodes compare the burst frequency at a time $s$ with the regulation function at the proteins one mean lifetime earlier, and the free-running fit reads the burst frequency of the simulated cells at the same lag. The benchmark of Section~\ref{sec:benchmark} draws the counts without memory and does not use it.

%======================================================================
\section{Families of regulation functions}\label{sec:families}
%======================================================================

\subsection{General form}

The burst frequency keeps the form~\eqref{eq:kon}, with
\begin{equation}\label{eq:family}
\sigma_i(P) = \ell_i\Bigl(\beta_i + \sum_j \theta_{ij}\,\varphi_i(P_j)\Bigr),
\end{equation}
a link $\ell_i$ (increasing, from $0$ to $1$) and a transformation $\varphi_i$ of the regulators (increasing), both chosen by target $i$. The network is $\theta$ for every family: $\theta_{ij} \neq 0$ means that $j$ regulates $i$, its sign gives activation or repression, and the sparsity penalty, the prior network and the refinement of CardamomOT apply unchanged. The family of a target is a discrete parameter $f_i = (\ell_i, \varphi_i) \in \F$.

\subsection{Sigmoid of a linear input}

$\ell = \sigma$, $\varphi(P) = P$: the current model. Each regulator multiplies the odds of the high mode by $e^{\theta_{ij}P_j}$, without saturation: a regulator that is twice as abundant has the square of its effect on the odds.

\subsection{Hill type}

$\ell = \sigma$, $\varphi(P) = \bar P\,\log(1 + P/\kappa)/\log(1 + \bar P/\kappa)$, with $\kappa > 0$ a scale of the regulator and $\bar P$ the protein level of its high mode:
\begin{equation}\label{eq:hill}
\sigma_i(P) = \frac{A_i\prod_j (1 + P_j/\kappa_j)^{c_j\theta_{ij}}}{1 + A_i\prod_j (1 + P_j/\kappa_j)^{c_j\theta_{ij}}}, \qquad A_i = e^{\beta_i},\quad c_j = \frac{\bar P_j}{\log(1 + \bar P_j/\kappa_j)}.
\end{equation}
The normalisation gives $\varphi(0) = 0$ and $\varphi(\bar P) = \bar P$, as the identity: in both families, $\theta_{ij}\bar P_j$ is the change of the logit of $i$ when $j$ goes from $0$ to its high mode. The interactions of the two families are thus in the same units, which matters for the inference: the sparsity penalty acts on $|\theta_{ij}|$, and without normalisation the Hill interactions, smaller by the factor $c_j$ for the same effect, would be penalised less, which biases the selection towards the Hill family.

\begin{prop}\label{prop:hill}
For a single regulator $j$ and $P_j \gg \kappa_j$, \eqref{eq:hill} is a Hill function of $P_j$ with coefficient $n = c_j\theta_{ij}$ and dissociation constant $K = \kappa_j A_i^{-1/n}$: $\sigma_i \approx P_j^n / (K^n + P_j^n)$. For $\kappa_j \to 0$ with $\beta_i + n\log(1/\kappa_j)$ fixed, the approximation is exact.
\end{prop}

The family has the same parameters as the sigmoid: the Hill coefficients are the interactions (up to the known factor $c_j$), and the dissociation constants are carried by the basal. $\kappa$ plays the role of the offset $\varepsilon$ of $\log(\varepsilon + P)$: it avoids $\log 0$, sets the leak of the function at $P = 0$, and the constant $-\theta_{ij}\log\kappa_j$ is absorbed by $\beta_i$. It is not inferred: $\kappa_j = \kappa\,\bar P_j$, with $\bar P_j$ the protein level of the high mode of $j$ (known from the mixture; $\bar P_j = 1$ in the normalised model) and $\kappa = 0.1$ (\texttt{model\_learning\_kappa}). With several regulators, \eqref{eq:hill} multiplies their contributions to the odds, as the sigmoid does, but each contribution grows as a power of the regulator instead of an exponential: the function is concave in $P_j$ and saturates in each regulator only through the link.

\subsection{The thermodynamic model of Herbach et al.}\label{sec:herbach}

The model of~\cite{herbachinferring2017}, from which CardamomOT derives, defines
\[
\Phi_i(P) = e^{\theta_{ii}}\prod_{j \neq i}\frac{1 + e^{\theta_{ij}}(P_j/s_{ij})^{m_{ij}}}{1 + (P_j/s_{ij})^{m_{ij}}}, \qquad
\koni(P) = \frac{k_{0,i} + k_{1,i}\,\Phi_i(P)\,(P_i/s_{ii})^{m_{ii}}}{1 + \Phi_i(P)\,(P_i/s_{ii})^{m_{ii}}}.
\]

\begin{prop}\label{prop:herbach}
The model of~\cite{herbachinferring2017} is of the form~\eqref{eq:kon}--\eqref{eq:family} with the logistic link:
\[
\koni(P) = k_{0,i} + (k_{1,i} - k_{0,i})\,\sigma\Bigl(\theta_{ii} + m_{ii}\log\frac{P_i}{s_{ii}} + \sum_{j\neq i}\psi(\theta_{ij}, P_j/s_{ij})\Bigr),
\qquad \psi(\theta, x) = \log\frac{1 + e^{\theta}x^m}{1 + x^m}.
\]
\end{prop}

\begin{proof}
$\koni = k_0 + (k_1 - k_0)\,u/(1 + u)$ with $u = \Phi_i(P)(P_i/s_{ii})^{m_{ii}}$, and $u/(1 + u) = \sigma(\log u)$.
\end{proof}

The self-regulation is a Hill term of fixed coefficient $m_{ii}$, the $\varphi = \log$ family of the previous section, and each other regulator enters through a \emph{saturating} contribution $\psi$, which goes from $0$ at $P_j = 0$ to $\theta_{ij}$ when $P_j \to \infty$, with a Hill transition of coefficient $m_{ij}$ around $s_{ij}$. Biologically it is the most realistic of the three: a transcription factor that saturates its binding sites has no further effect. The cost is the parameters $s_{ij}$ and $m_{ij}$, and $\psi$ is nonlinear in $\theta$.

\paragraph{Approximate form.} With $h_m(x) = x^m/(1 + x^m)$ the fraction of occupied sites, $\psi(\theta, x) = \log\bigl(1 + (e^{\theta} - 1)h_m(x)\bigr) = \theta\,h_m(x) + O(\theta^2)$: to first order in $\theta$, each regulator contributes $\theta$ times a Hill step. The approximate Herbach family keeps this first order, normalised as the Hill type,
\begin{equation}\label{eq:herbach-approx}
\varphi(P_j) = \bar P_j\,\frac{h_m(P_j/s_j)}{h_m(\bar P_j/s_j)}, \qquad \sigma_i(P) = \sigma\Bigl(\beta_i + \sum_j \theta_{ij}\,\varphi(P_j)\Bigr),
\end{equation}
so that $\varphi(0) = 0$ and $\varphi(\bar P_j) = \bar P_j$: as for the identity and the Hill type, $\theta_{ij}\bar P_j$ is the change of the logit of $i$ when $j$ goes from $0$ to its high mode, and the three families share their endpoints and the units of $\theta$. It is the form used for now (\texttt{model\_learning\_herbach\_form = 'approximate'}). Being linear in $\theta$, it is fitted exactly as the other families, on the features $\varphi(P_j)$ of the regulators, with the same loss, sparsity threshold and ridge. It differs from the exact form when $|\theta_{ij}|$ is large: the exact $\psi$ is concave in the occupancy $h_m$ for an activator and convex for a repressor, while~\eqref{eq:herbach-approx} is linear in it. The self-regulation is treated as any other regulator. The exact form (\texttt{'exact'}, nonlinear fit) is kept for later study.

\paragraph{Threshold and Hill coefficient.} The threshold $s_j$ is a property of the regulator, not inferred: halfway between its low and high protein levels, $s_j = \bar P_j(1 + k_{0,j}/k_{1,j})/2$, from the lowest and highest modes of its mixture ($\bar P_j/2$ for a stimulus). The Hill coefficient $m$ is a property of the \emph{target}, shared by its regulators (the cooperativity of its promoter), and takes a few discrete values, $m \in \{1, 2, 4\}$ (\texttt{model\_learning\_herbach\_m}). Each value is a candidate family of its own, so that $m$ is selected with the family, by the same held-out score (Section~\ref{sec:selection}), and the family has no continuous parameter beyond $\theta$ and $\beta$. With $m = 1$, $\varphi$ is a concave saturation (Michaelis--Menten of half-saturation $s_j \approx \bar P_j/2$), which matches the Hill type of the previous section ($\kappa = 0.1$) on $[0, \bar P_j]$ within $0.04$--$0.08$: the two families differ only beyond the high mode, where the Hill type keeps growing as a logarithm, and the Hill type is therefore not a default candidate (two indistinguishable candidates would share the targets at random); with $m = 2$, $\varphi$ is a smooth sigmoidal step, within $0.09$--$0.14$ of the identity; with $m = 4$, it is nearly a step at $s_j$, and the target switches when its regulator crosses the threshold, whatever its level beyond. The coefficients differ only by the steepness of $\varphi$ around $s_j$, hence only on the states whose regulators are near their threshold, i.e.\ along the transients.

\subsection{Richards link}\label{sec:richards}

The generalised logistic of Richards~\cite{richards1959flexible}, $\ell_\nu(z) = (1 + e^{-z})^{-1/\nu}$, $\nu > 0$, is the logistic for $\nu = 1$. For $\nu \neq 1$ it is asymmetric: its inflection is at $\ell = (1 + \nu)^{-1/\nu} \neq 1/2$, so the switch towards the high mode and the switch back have different steepness (a gene that needs a large excess of activator to switch on, and switches off quickly, or the opposite). It adds one parameter $\nu_i$ per gene, and changes the shape of the response, not the way regulators combine.

\subsection{Co-regulation within one network}\label{sec:coreg}

The families above combine the regulators independently on the odds: $\log$-odds $= \beta_i + \sum_j \theta_{ij}\varphi_i(P_j)$. Synergy (two activators required together) or antagonism (one regulator cancelling another) needs interaction terms. Products of per-regulator functions (AND / OR gates) are the classical answer in Boolean and thermodynamic models; in the present setting, the factors would be occupancies $f_{ij}(P_j)$ of the binding sites of the regulators on the promoter of $i$, functions of the protein levels, not states of the regulator genes. A lighter version that keeps one network is a pairwise term
\[
\beta_i + \sum_j \theta_{ij}\varphi_i(P_j) + \sum_{j<k}\gamma_{ijk}\,\varphi_i(P_j)\varphi_i(P_k),
\]
with $\gamma_{ijk}$ defined only for pairs of regulators already in the support of $\theta_{i\cdot}$, and a strong $L^1$ penalty: $\gamma_{ijk} > 0$ is a co-activation, $\gamma_{ijk} < 0$ a redundancy. The network is still $\theta$, and $\gamma$ is a second-order annotation of its edges. The thermodynamic models of~\cite{shea1985lambda, bintu2005transcriptional} contain both forms as limits, at the price of more parameters.

%======================================================================
\section{Identifiability of the family}\label{sec:identifiability}
%======================================================================

\begin{prop}\label{prop:steady}
Let target $i$ have a single regulator $j$, and let the observed states have $P_j \in \{0, \bar P_j\}$ only (no transient). Any family of~\eqref{eq:family} with increasing $\ell$ and $\varphi$ fits the basins of $i$ exactly when $\theta_{ij}$ and $\beta_i$ can be chosen so that $\sigma_i(0) = \tau_i^{\mathrm{low}}$ and $\sigma_i(\bar P_j) = \tau_i^{\mathrm{high}}$: the family is not identifiable.
\end{prop}

Two points fix two parameters. The family is identified by states at intermediate protein levels, and these exist along transients: after a switch of the regulator, $P_j$ relaxes at rate $d_{1,j}$, and the target switches when $\sigma_i(P_j(t))$ crosses the level that its basins indicate. A sigmoid of a linear input switches near the middle of the range of $P_j$ (for a symmetric $\beta_i$), a Hill target of the same endpoints switches earlier and saturates later, because $\log(1 + P/\kappa)$ is concave. The information is in the \emph{timing} of the switches relative to the protein levels, which the mechanistic bridge reconstructs; it grows with the protein lifetime (slow transients are sampled by more timepoints) and with the number of cells caught during the switch. The benchmark of Section~\ref{sec:benchmark} is built so that both families share their endpoints, and only this information separates them. Whether it survives the reading of the targets on the counts, then the reconstruction of the proteins, is the question of the benchmark (Section~\ref{sec:benchmark}).

%======================================================================
\section{Selection of the family in the inference loop}\label{sec:selection}
%======================================================================

\paragraph{Criterion.} At an update of the network, the network is fitted with every candidate family $f$ for every target, on the same sub-samples of trajectory states as the standard fit (\texttt{n\_network\_fits} disjoint sub-samples, stratified by time, sample and cell type). The fit $\hat\theta^{(f,k)}$ on sub-sample $k$ is scored on the states $\bar{\mathcal{S}}_k$ of the other sub-samples by the likelihood of the counts of their real cells under the network prior on the levels (\texttt{model\_learning\_score = 'marginal'}),
\begin{equation}\label{eq:score}
S_i(f) = -\frac{1}{K}\sum_{k=1}^K \frac{1}{|\bar{\mathcal{S}}_k|}\sum_{n \in \bar{\mathcal{S}}_k} \log \sum_l \mathrm{NB}(x_{ni};\, \mathbf{k}_{l,i})\, p\bigl(l \mid \bar\sigma_i^{(f,\hat\theta^{(f,k)})}(n)\bigr), \qquad p(l \mid \sigma) \propto e^{-\kappa_i \mathrm{KL}(\tau_{l,i} \| \sigma)},
\end{equation}
with $\bar\sigma_i^{(f,\theta)}(n)$ the regulation function of family $f$ along the memory of the mRNA~\eqref{eq:memory} and $\kappa_i$ the precision of the basin update. It is the likelihood~\eqref{eq:complete} with the basin summed out, the quantity maximised by the EM; a cell whose counts do not discriminate the levels contributes the same constant to every family. $\mathrm{NB}(x_{ni};\mathbf{k}_{l,i})$ is read as $\pi_{ni}(l)/w_{l,i}(t_n)$, the posterior of the mixture over the mass of the level at the time of the cell, up to a constant per cell. The held-out states let families with more parameters be compared without an explicit penalty. The families are compared at their unshrunk effects: the refinement of the network fit re-estimates the effects on the support and signs of the penalised fit without the $L^1$ penalty, keeping its ridge, which bounds the edges of nearly separable targets (relaxed lasso, \texttt{network\_relaxed}); with shrunk effects, a sigmoid target is better approximated by a saturating function than by a shrunk sigmoid, and goes to the Herbach type.

The alternative score (\texttt{'ce'}) is the cross-entropy against the expected levels $\bar\tau_{ni} = \E_q[\tau_{z_{ni},i}]$ under the current posteriors, the $Q$ function of the M-step (Proposition~\ref{prop:mstep}). It favours the current family $f_0$: a cell whose counts do not discriminate the levels has $\bar\tau_{ni} \approx \bar\sigma_i^{(f_0)}(n)$, a cost $0$ for $f_0$ and $\mathrm{KL}(\bar\sigma^{(f_0)} \,\|\, \bar\sigma^{(f)}) > 0$ for the others, so that only the informative cells move a target to another family.

\paragraph{Candidates.} The candidates are the products of the transformations and links, the Herbach type once per Hill coefficient: with the defaults, identity and Herbach with $m = 1, 2, 4$, four families of two types. The Hill coefficient is thus chosen as a family, by the same score, and never fitted continuously.

\paragraph{Starting family.} Every target starts as a sigmoid of a linear input, the family of two modes, and leaves it when the score~\eqref{eq:score} of another family is lower. The weight of the intermediate levels of the mixture, a prior on the family, recovers only $0.28$ of the types on FN4 and is not used.
\paragraph{Schedule.} Before the selection window (iterations $< n_{\mathrm{start}}$, by default \texttt{min\_n\_loops}), the trajectories, the transport and the basins stabilise as in CardamomOT, while the weight of the network in the costs~\eqref{eq:costs} and the transport are ramped up; every target is a sigmoid, and the targets of the network fit ramp from the EMD basins to the expected levels, so that the first selection is made on settled posteriors. From $n_{\mathrm{start}}$ on, at every iteration (\texttt{model\_learning\_period} $= 1$), each update runs all the candidate fits; target $i$ moves to the family $\argmin_f S_i(f)$, type and $m$ together, when its score is lower than that of the current family by a relative margin (\texttt{model\_learning\_margin}) at two consecutive selections, i.e.\ two consecutive iterations (hysteresis against the noise of the sub-samples and against oscillations, since the next trajectories are built with the selected function). The parameters of each target are those of the fit of its family. Each family keeps its own $\theta$ from one update to the next (warm start); a family fitted for the first time starts from the current $\theta$, which has the same units for every $\varphi$. The families are frozen from $n_{\mathrm{freeze}} = 2\,n_{\mathrm{start}}$ on (\texttt{model\_learning\_freeze}), after which the loop converges with them, so that its stopping criterion keeps its meaning.


\paragraph{As an EM.} The family index, Hill coefficient included, is a discrete parameter of the M-step of Section~\ref{sec:em}: with the posteriors (E-step) and the trajectories fixed, each target maximises over $(f, \theta_{i\cdot}, \beta_i)$, the continuous parameters by the network fit (Proposition~\ref{prop:mstep}) and the family by the held-out likelihood. After the freeze, the procedure is the EM of Section~\ref{sec:em} on a fixed model.

\paragraph{Cost.} The network fits of the targets are independent and parallel: a selection multiplies the cost of the network update by $|\F|$ (four with the defaults), at each iteration of the window $[n_{\mathrm{start}}, 2n_{\mathrm{start}})$. With a single candidate family, it is imposed on every target from the first iteration, without selection: the whole pipeline (inference, kinetic rates, simulations) then runs with that regulation function, at the cost of the standard method.

\paragraph{After the loop.} The kinetic rates are recalibrated by the free-running fit of the simulation model (refined $\theta$, rates and scale of $\theta$), which uses the selected families through the torch mirror of the regulation function. The generative loss of this fit, with all genes sigmoid against the selected families, is a validation of the selection at the level of the whole model.

%======================================================================
\section{Implementation}\label{sec:implementation}
%======================================================================

\begin{itemize}[nosep]
\item \texttt{CardamomOT/inference/sigma.py}: specification \texttt{spec} $(G, 6)$: $\varphi$, link and Hill coefficient $m$ (Herbach) of each gene as a target, scale $\kappa$, high-mode level $\bar P$ and threshold $s$ (Herbach) of each gene as a regulator; kernels \texttt{sigma\_pairs}, \texttt{kon\_matrix}, \texttt{kon\_last}, \texttt{kon\_upper\_bound} (thinning bound of the explicit-mRNA PDMP, valid because $\varphi$ and $\ell$ are increasing), and \texttt{sigma\_torch}. \texttt{spec = None} reproduces the sigmoid of a linear input bit for bit.
\item Mixture (\texttt{inference/mixture.py}): classification EM and soft EM with masses per timepoint (\texttt{max\_iter\_softEM}); continuous model~\eqref{eq:grid} (\texttt{fit\_grid\_mixture}; \texttt{max\_grid\_levels}, \texttt{mixture\_selection\_criterion}, \texttt{mixture\_criterion}).
\item Loop (\texttt{model/base.py}, \texttt{fit\_network}): costs~\eqref{eq:costs} (\texttt{\_mode\_network\_cost}, precision $\kappa_i$, \texttt{grid\_network\_precision\_max}), EMD basins and entropic posteriors (\texttt{\_assign\_from\_costs}, \texttt{one\_hot\_probas}), memory of the mRNA~\eqref{eq:memory} (\texttt{targets\_accounting\_late\_mrna}, \texttt{late\_mrna\_nodes}).
\item Network fit (\texttt{inference/network.py}): a target with $\varphi \neq$ identity is fitted on the features $\varphi(P_j)$ of its regulators, with the same loss, penalties and refinement; $\bar\sigma$ on the memory nodes.
\item Selection (\texttt{model/base.py}): \texttt{\_select\_sigma\_families}, called by the network update of \texttt{fit\_network} during the window; parameters \texttt{model\_learning\_grn}, \texttt{model\_learning\_phi}, \texttt{model\_learning\_link}, \texttt{model\_learning\_herbach\_form} (\texttt{'approximate'}), \texttt{model\_learning\_herbach\_m} (\texttt{[1, 2, 4]}), \texttt{model\_learning\_kappa}, \texttt{model\_learning\_start}, \texttt{model\_learning\_period}, \texttt{model\_learning\_freeze}, \texttt{model\_learning\_margin}, \texttt{model\_learning\_score} (\texttt{'marginal'}, \texttt{'ce'}), \texttt{network\_relaxed}; thresholds $s_j$ by \texttt{\_herbach\_thresholds}.
\item Observed proteins (\texttt{infer\_network\_structure.py}, \texttt{Data/data\_prot.h5ad}, same cells and genes as the counts, any unit per gene): \texttt{model.observed\_prot}; \texttt{fit\_network} then calls \texttt{fit\_observed\_proteins} instead of \texttt{loop\_trajectories}: scale of the proteins (\texttt{\_protein\_scale}, saved in \texttt{protein\_scale.npy}), weights without branching (\texttt{\_unbranched\_log\_weights}), EM on the real cells (network fit, then \texttt{\_update\_basins}), and states of the later steps resampled with the weights. Without the file, nothing changes.
\item Outputs: \texttt{cardamomOT/sigma\_spec\_grn.npy} (used by every later step: simulations, test set, report) and \texttt{sigma\_selection.npz} (scores~\eqref{eq:score} and families at each selection).
\end{itemize}

%======================================================================
\section{Benchmark}\label{sec:benchmark}
%======================================================================

The benchmark separates the difficulties. The data are simulated with the model of the inference itself, without the mRNA memory and the explicit mRNA dynamics of the full model, and the selection is evaluated first with the proteins observed, then hidden.

\paragraph{Networks.} FN4, CN5, BN8 and FN8, the simulated networks of the CardamomOT article, with their degradation rates, basals, interactions and the parameters of Harissa ($k_0 = 0$, $k_1 = 2$ in units of $d_0$, mRNA burst size $50$). Each gene is, with probability $1/4$ drawn per run, a sigmoid target or a Herbach target in its approximate form (Section~\ref{sec:herbach}; threshold $1/2$) of Hill coefficient $m = 1$, $2$ or $4$. With the normalised $\varphi$, every family keeps the interactions of the network, has the same logit at $P = 0$ and at the high mode, and switches between the same burst frequencies with a different shape: only the transients separate them (Section~\ref{sec:identifiability}).

\paragraph{Simulation.} Model of the inference: bursty PDMP of the proteins only (\texttt{simulate\_next\_prot\_pdmp}, parametrised as the simulations of CardamomOT, high mode at $P = 1$), independent cells with a burn-in of $5$\,h without stimulus, stimulus from $t = 0$, $1000$ cells over the ten timepoints $0, 6, \dots, 96$\,h. The counts of a cell observed at time $t$ are drawn from the negative binomial of its burst frequency at $t$, $\mathrm{NB}\bigl(\koni(P(t))/d_{0,i},\, c_i\bigr)$: the observation model of the mixture, without memory of the mRNA. Ten runs per network (\texttt{projet\_modellearning/generate.py}).

\paragraph{Three levels.}
\begin{enumerate}[nosep]
\setcounter{enumi}{-1}
\item Oracle (\texttt{oracle\_families.py}): the network fit and the held-out score of the selection on the true proteins of the cells, without the loop, against the true burst frequency (noise-free upper bound) or against the targets of the mixture.
\item Observed proteins: the proteins of the cells are given to CardamomOT (\texttt{Data/data\_prot.h5ad}), not their burst frequencies, in an unknown unit per gene. The mixture is unchanged. Each cell then carries both its regulators $P$ and the posterior of its target levels, so no trajectory is needed: the inference is a plain EM on the real cells, the network fitted on $(P, \text{expected levels})$ (with the selection of the families) and the basins updated given the network, without trajectories nor switching times $\alpha$. This is the ideal case of the model learning: the true model of the inference, known proteins, and targets read on the counts only. Experimentally, it is approached for a few proteins by CITE-seq.

\emph{Scale of the proteins.} The observed proteins are $\lambda_g P_g$, with $P_g$ in the units of the model (deterministic high mode at 1). Since the protein decay is linear, the first moment of the PDMP is exact whatever the network: $\frac{d}{dt}\mathbb{E}[P_g] = d_{1,g}\,\mathbb{E}[u_g] - \mathbb{E}[(d_{1,g} + b)P_g]$, with $u_g = k_{\mathrm{on},g}/k_{\mathrm{high},g}$ and $b$ the birth rate (dilution). $\mathbb{E}[u_g](t)$ is read on the mRNA alone (expected modes under the posteriors of the mixture), and $\mathbb{E}[bP_g]/\mathbb{E}[P_g]$ is free of the unit. The observed means thus follow $\lambda_g m_g(t)$, $\dot m_g = d_{1,g}\bar u_g(t) - (d_{1,g} + \bar b_g(t)) m_g$, and $\lambda_g$ (with the initial state of each sample, from three timepoints) is fitted by non-negative least squares. The proteins are divided by $\lambda_g$: a change of unit, which keeps their distribution.

\emph{Proliferation.} The observed cells sample the branched population, the model describes the SDE without branching: each cell is weighted by $w = \mathbb{E}[e^{-\int R}]$ along its lineage, the inverse of the branching of the simulations, by recursion between consecutive timepoints with ancestors drawn by a Gaussian kernel on the proteins, $p(i\mid j) \propto e^{R_i\Delta t/2 - |P_i - P_j|^2/2h^2}$, and $w_j = \sum_i p(i\mid j)\, w_i\, e^{-(R_i + R_j)\Delta t/2}$. The scale is fitted without the weights, then with them, and the final weights (on the rescaled proteins) resample the states of the later steps (refine, simulations). The benchmark has no proliferation (weights 1).
\item Hidden proteins: the same data, standard CardamomOT (proteins reconstructed from the modes read on the mRNA), with the options of the simulated model (no mRNA bridge, no memory). The difference with level 1 measures the cost of the reconstruction of the proteins: these are known up to a monotone distortion $\hat P_j = h_j(P_j)$ specific to each regulator, and the family fitted on them is $\varphi\circ h_j^{-1}$, not $\varphi$.
\end{enumerate}

\paragraph{Metrics.} Accuracy of the family of each gene, $m$ included (chance level $1/4$ for a uniform choice among the four candidates), accuracy of its type (sigmoid or Herbach, whatever $m$; chance $5/8$ with the uniform draw of the families), recall of each type, and accuracy of the Hill coefficient of the Herbach targets classified as Herbach (chance $1/3$); AUPR of the inferred network (self-loops and edges towards the stimulus excluded), sign of the true edges and mean absolute error on $\theta$, with and without model learning (\texttt{run\_inference.py}, modes \texttt{obs\_ml}, \texttt{obs\_base}, \texttt{hid\_ml}, \texttt{hid\_base}).

%======================================================================
\section{Results}\label{sec:results}
%======================================================================

\tofill{levels 0, 1 and 2 on the four networks, ten runs each}

%======================================================================
\section{Open questions}\label{sec:open}
%======================================================================

\begin{itemize}[nosep]
\item Slower proteins, which keep the regulators longer at intermediate levels, and more contrasted families (smaller $\kappa$).
\item The exact distribution of the counts given the history of the burst frequency (shot noise with memory), which would make the likelihood of the counts a valid score.
\item Soft selection: weights $w_{if} \propto e^{-N S_i(f)}$ instead of a hard choice, or a mixture of families, at the price of identifiability between the $\theta$ of the families.
\item Selection by the generative loss of the free-running fit, gene by gene, which tests the families on the simulated dynamics and not only on the teacher-forced trajectories.
\item Families shared by groups of genes (same transcription factor binding mode) instead of one per target.
\end{itemize}

%======================================================================
\begin{thebibliography}{99}
%======================================================================

\bibitem{mauge2026cardamomot}
Y.~Maugé, E.~Ventre.
\newblock CardamomOT: a mechanistic optimal transport-based framework for gene regulatory network inference, trajectory reconstruction and generative modeling.
\newblock \emph{bioRxiv}, 2026. doi:10.64898/2026.03.31.715390. Accepted in \emph{PLoS Computational Biology}.

\bibitem{herbachinferring2017}
U.~Herbach, A.~Bonnaffoux, T.~Espinasse, O.~Gandrillon.
\newblock Inferring gene regulatory networks from single-cell data: a mechanistic approach.
\newblock \emph{BMC Systems Biology}, 11:105, 2017.

\bibitem{shea1985lambda}
M.~A.~Shea, G.~K.~Ackers.
\newblock The $O_R$ control system of bacteriophage lambda: a physical-chemical model for gene regulation.
\newblock \emph{Journal of Molecular Biology}, 181(2):211--230, 1985.

\bibitem{bintu2005transcriptional}
L.~Bintu, N.~E.~Buchler, H.~G.~Garcia, U.~Gerland, T.~Hwa, J.~Kondev, R.~Phillips.
\newblock Transcriptional regulation by the numbers: models.
\newblock \emph{Current Opinion in Genetics \& Development}, 15(2):116--124, 2005.

\bibitem{brunton2016sindy}
S.~L.~Brunton, J.~L.~Proctor, J.~N.~Kutz.
\newblock Discovering governing equations from data by sparse identification of nonlinear dynamical systems.
\newblock \emph{Proceedings of the National Academy of Sciences}, 113(15):3932--3937, 2016.

\bibitem{mangan2016inferring}
N.~M.~Mangan, S.~L.~Brunton, J.~L.~Proctor, J.~N.~Kutz.
\newblock Inferring biological networks by sparse identification of nonlinear dynamics.
\newblock \emph{IEEE Transactions on Molecular, Biological and Multi-Scale Communications}, 2(1):52--63, 2016.

\bibitem{richards1959flexible}
F.~J.~Richards.
\newblock A flexible growth function for empirical use.
\newblock \emph{Journal of Experimental Botany}, 10(2):290--301, 1959.

\bibitem{delyon1999saem}
B.~Delyon, M.~Lavielle, E.~Moulines.
\newblock Convergence of a stochastic approximation version of the EM algorithm.
\newblock \emph{Annals of Statistics}, 27(1):94--128, 1999.

\bibitem{shen1998triple}
W.~Shen, T.~A.~Louis.
\newblock Triple-goal estimates in two-stage hierarchical models.
\newblock \emph{Journal of the Royal Statistical Society: Series B}, 60(2):455--471, 1998.

\end{thebibliography}

\end{document}
